\chapter*{Abstract}
\addcontentsline{toc}{chapter}{Abstract}

This thesis addresses weakly supervised mutation prediction from
gigapixel H\&E (haematoxylin-and-eosin-stained) images of lung
adenocarcinoma (LUAD, the most common subtype of lung cancer): a
single binary label per gene is assigned to $\sim\!10^9$ pixels
with no sub-region annotation. Existing systems require tens of thousands
of proprietary slides and offer no explanation of \emph{why} a
mutation is predicted---a critical barrier to clinical trust. The
underlying problem is itself \emph{stochastic, many-to-many and
weakly labelled}: the same morphology can correspond to multiple
genotypes and vice versa, so competence must be argued per gene
and per slide rather than from a single cohort metric.

The central hypothesis is that \emph{structured domain
knowledge---specifically, the six expert-defined LUAD histological
growth patterns
(lepidic, acinar, papillary, micropapillary, solid, cribriform)
encoded through fuzzy set theory---can be injected into a deep
learning pipeline to simultaneously achieve competitive mutation
prediction under data scarcity ($<\!700$ public slides) and render
the resulting predictions \emph{interpretable} through a
multi-level, clinically verifiable inspection framework---and,
through that framework, \emph{explainable} via
knowledge-graph--grounded natural-language rationales---to a
degree that no end-to-end black-box approach offers}; the verdict
is \emph{gene-conditional}.
Following design science research, three artefacts and one
integrative system (\lchai{}, \emph{Lung Cancer Histologic
Analysis with AI}) are designed and evaluated.

\textbf{Artefact~1 (FuzzyArcLoss~V2)} is a sample-wise fuzzy
modulation of the ArcFace angular margin that respects gradual
morphological transitions between pattern classes; it ranks
first among 18 losses (macro-F1 $= 92.31\%\!\pm\!2.04$,
$p\!=\!0.001$, Cohen's $d\!=\!1.00$). Its advantage is
\emph{morphological-boundary-dependent, not class-size-dependent}:
the gain materialises on visually ambiguous adjacent classes
irrespective of class size---a transfer rule that generalises to
Gleason grading, Pap-smear cytology and dermatopathology.

\textbf{Artefact~2 (PI-ABMIL)} concatenates a 512-d CTransPath
embedding with the 6-d fuzzy pattern probabilities from
Artefact~1 before gated-attention aggregation, anchoring
slide-level reasoning in clinically validated patterns.

\textbf{Artefact~3 (FC-MIL)} replaces additive attention with a
learned 2-additive Choquet integral whose 21 parameters (six
Shapley values, fifteen pairwise interactions) capture
supra-additive interactions between morphological categories
that standard attention cannot represent.

Evaluation on TCGA-LUAD (\emph{The Cancer Genome Atlas Lung
Adenocarcinoma}; 687 slides, 668 patients, 5-fold
patient-stratified CV) over six driver mutations
(\textit{TP53}, \textit{EGFR}, \textit{KRAS}, \textit{STK11},
\textit{KEAP1}, \textit{RBM10}; descending prevalence) shows
pattern-knowledge utility is \emph{gene- and
aggregator-conditional}, decomposing into three structural
regimes:

\textbf{Regime~(i) --- encoder saturation.}
The pretrained encoder already captures the relevant morphology
and the pattern channel adds nothing
(\textit{TP53}, \textit{EGFR}).

\textbf{Regime~(ii) --- inter-pattern co-occurrence.}
The discriminative signal lives in pairwise pattern interactions
rather than in any single pattern; only an aggregator that reads
joint activations recovers it
(\textit{KRAS}, via FC-MIL's Choquet integral).

\textbf{Regime~(iii) --- out-of-ontology.}
No representation within the 6-pattern simplex
(ANORAK~\cite{anorak2024}) carries the signal; the regime
decomposes into three named mechanisms---\emph{label-space
mismatch}~\cite{lipton2018labelshift}, \emph{no morphological
projection}, and
\emph{collapse-to-prior}~\cite{bowman2016posteriorcollapse}---that
together bound what H\&E-only mutation prediction can achieve
(\textit{STK11}, \textit{KEAP1}, \textit{RBM10}).

To expose these limits, \lchai{} instruments a six-level
interpretability stack (attention, pattern overlays, SHAP,
Choquet Shapley values, channel ablation, fuzzy linguistic
labels) and an LLM explanation agent grounded in a curated
knowledge graph built from NCIt and MONDO ontologies: fuzzy
labels constrain \emph{how} the LLM speaks, the graph
constrains \emph{what} it can claim---a dual anti-hallucination
mechanism.
A complementary cohort-level Conclusive/Inconclusive flag
(derived from AUROC) gates the per-slide explanation pathway,
separating population-level signal presence from per-slide
morphological confidence.

The hypothesis is \emph{partially supported}: pattern knowledge
yields interpretability that no prior mutation-prediction system
offers, and a genuine inter-pattern predictive gain on KRAS; for
the remaining genes the bound lies in the 6-pattern \emph{input
representation} itself (the ANORAK label space), not in the
downstream model
(loss\,/\,feature fusion\,/\,aggregator).
The principal claims of this thesis are therefore not
clinical-grade accuracy---foundation-scale 60{,}000+ slide
cohorts already exist for that---but the following two,
distilled from the six fine-grained contributions C1--C6
listed in Chapter~\ref{ch:introduction}:

\textbf{Claim~(i) --- methodological.}
Fuzzy logic, applied coherently across loss, representation and
aggregation, makes a public-data pipeline transparent through a
knowledge-graph--grounded six-level interpretability stack.

\textbf{Claim~(ii) --- empirical/analytical.}
The gene-by-gene taxonomy derived from per-slide ablation on
TCGA-LUAD bounds what any morphology-only model can achieve from
H\&E alone, and identifies the genes for which fusing additional
data sources (e.g., gene-expression profiles or radiology),
rather than further architectural refinement, is the natural
next step.

Together, these claims position the approach as a complementary
paradigm to foundation-scale brute-force---one in which a
pathologist can understand, verify and \emph{correct} the model's
reasoning, and in which interpretability is most valuable
precisely on the cases the predictor gets wrong.

\paragraph{Keywords.}
Fuzzy logic, angular margin loss, multiple-instance learning,
Choquet integral, Shapley values, histopathological patterns,
lung adenocarcinoma, mutation prediction, genotype--phenotype gap,
design science research, computational pathology, explainable~AI.
